
\documentclass{article}
\usepackage{colortbl}
\usepackage{makecell}
\usepackage{multirow}
\usepackage{supertabular}

\begin{document}

\newcounter{utterance}

\centering \large Interaction Transcript for game `cladder', experiment `full\_v1.5\_default', episode 8085 with reflexion{-}llama{-}3.1{-}8b{-}instruct.
\vspace{24pt}

{ \footnotesize  \setcounter{utterance}{1}
\setlength{\tabcolsep}{0pt}
\begin{supertabular}{c@{$\;$}|p{.15\linewidth}@{}p{.15\linewidth}p{.15\linewidth}p{.15\linewidth}p{.15\linewidth}p{.15\linewidth}}
    \# & \multicolumn{2}{c}{Player} && \multicolumn{2}{c}{Game Master} \\
    \hline

    \theutterance \stepcounter{utterance}  
    & & & \multicolumn{4}{p{0.6\linewidth}}{
        \cellcolor[rgb]{0.9,0.9,0.9}{
            \makecell[{{p{\linewidth}}}]{
                \texttt{\tiny{[P1$\langle$GM]}}
                \texttt{You have attempted this same task before and did not succeed. These are the plans you wrote after those attempts.} \\
\texttt{Trial \#1: New plan: } \\
\\ 
\texttt{In my previous attempt, I failed to address the counterfactual scenario correctly. I considered the probability of kwoz for those who are kwox but not kwox as if it were a straightforward conditional probability. However, the question asks us to compare the probability of kwoz for those who are kwox but not kwox, which implies a counterfactual scenario.} \\
\\ 
\texttt{To address this correctly, I should have considered the potential outcomes under the counterfactual scenario. If an individual is kwox, the probability of kwoz is 65\%. If an individual is not kwox, the probability of kwoz is 35\%. However, the presence of kwox directly influences the probability of kwoz.} \\
\\ 
\texttt{To reason about the counterfactual scenario, I should have used the law of total probability to calculate the probability of kwoz for those who are kwox but not kwox. However, I should have also considered the effect of kwox on the probability of kwoz.} \\
\\ 
\texttt{A more accurate approach would be to use the concept of "backdoor adjustment" to reason about the counterfactual scenario. This involves adjusting the probability of kwoz for those who are kwox but not kwox by considering the effect of kwox on the probability of kwoz.} \\
\\ 
\texttt{In this case, I should have calculated the probability of kwoz for those who are kwox but not kwox as follows:} \\
\\ 
\texttt{P(kwoz|kwox, not kwox) = P(kwoz|not kwox) + P(kwox|not kwox) * (P(kwoz|kwox) {-} P(kwoz|not kwox))} \\
\\ 
\texttt{This formula takes into account the effect of kwox on the probability of kwoz and adjusts the probability of kwoz for those who are kwox but not kwox accordingly.} \\
\\ 
\texttt{Therefore, my new plan is to use the concept of backdoor adjustment to reason about the counterfactual scenario and calculate the probability of kwoz for those who are kwox but not kwox.} \\
\\ 
\\ 
\texttt{You are an expert in causal inference. The following question is not a typical commonsense query, but rather a meticulously designed question created by a professor specializing in causal inference, intended to assess the students' mastery of the course content.} \\
\texttt{Imagine a self{-}contained, hypothetical world with only the following conditions, and without any unmentioned factors or causal relationships: Kwox has a direct effect on kwoz. Swoq has a direct effect on kwoz. For those who are not kwox, the probability of kwoz is 35\%. For those who are kwox, the probability of kwoz is 65\%. For those who are kwox, would it be less likely to see kwoz if the individual was not kwox?} \\
\texttt{Start your answer with 'Yes' or 'No', followed by additional reasoning or evidence to support your explanation.} \\
            }
        }
    }
     \\ \\

    \theutterance \stepcounter{utterance}  
    & \multicolumn{4}{p{0.6\linewidth}}{
        \cellcolor[rgb]{0.561,0.737,0.561}{
            \makecell[{{p{\linewidth}}}]{
                \texttt{\tiny{[P1$\rangle$GM]}}
                \texttt{Yes.} \\
\\ 
\texttt{To address this question, we need to consider the counterfactual scenario where an individual is kwox but not kwox. This implies that we are comparing the probability of kwoz for those who are kwox but not kwox to the probability of kwoz for those who are not kwox.} \\
\\ 
\texttt{Using the concept of backdoor adjustment, we can calculate the probability of kwoz for those who are kwox but not kwox as follows:} \\
\\ 
\texttt{P(kwoz|kwox, not kwox) = P(kwoz|not kwox) + P(kwox|not kwox) * (P(kwoz|kwox) {-} P(kwoz|not kwox))} \\
\\ 
\texttt{However, we need to consider the effect of kwox on the probability of kwoz. Since kwox has a direct effect on kwoz, we need to adjust the probability of kwoz for those who are kwox but not kwox accordingly.} \\
\\ 
\texttt{To do this, we can use the formula:} \\
\\ 
\texttt{P(kwoz|kwox, not kwox) = P(kwoz|not kwox) + P(kwox|not kwox) * (P(kwoz|kwox) {-} P(kwoz|not kwox))} \\
\\ 
\texttt{We know that P(kwoz|not kwox) = 0.35, P(kwox|not kwox) = 0 (since an individual cannot be both kwox and not kwox), P(kwoz|kwox) = 0.65, and P(kwox|not kwox) = 0.} \\
\\ 
\texttt{However, we need to consider the effect of kwox on the probability of kwoz. Since kwox has a direct effect on kwoz, we can assume that the effect of kwox on kwoz is not zero.} \\
\\ 
\texttt{Let's denote the effect of kwox on kwoz as Δ. Then, we can rewrite the formula as:} \\
\\ 
\texttt{P(kwoz|kwox, not kwox) = P(kwoz|not kwox) + Δ} \\
\\ 
\texttt{Since P(kwoz|not kwox) = 0.35, we can rewrite the formula as:} \\
\\ 
\texttt{P(kwoz|kwox, not kwox) = 0.35 + Δ} \\
\\ 
\texttt{Now, we need to determine the value of Δ. Since kwox has a direct effect on kwoz, we can assume that the effect of kwox on kwoz is not zero.} \\
\\ 
\texttt{Let's assume that the effect of kwox on kwoz is Δ = 0.3 (this is a reasonable assumption, since the probability of kwoz for those who are kwox is 65\%,} \\
            }
        }
    }
    & & \\ \\

    \theutterance \stepcounter{utterance}  
    & & & \multicolumn{2}{p{0.3\linewidth}}{
        \cellcolor[rgb]{0.9,0.9,0.9}{
            \makecell[{{p{\linewidth}}}]{
                \texttt{\tiny{[GM$|$GM]}}
                \texttt{Yes} \\
            }
        }
    }
    & & \\ \\

    \theutterance \stepcounter{utterance}  
    & & & \multicolumn{2}{p{0.3\linewidth}}{
        \cellcolor[rgb]{0.9,0.9,0.9}{
            \makecell[{{p{\linewidth}}}]{
                \texttt{\tiny{[GM$|$GM]}}
                \texttt{yes} \\
            }
        }
    }
    & & \\ \\

    \theutterance \stepcounter{utterance}  
    & & & \multicolumn{2}{p{0.3\linewidth}}{
        \cellcolor[rgb]{0.9,0.9,0.9}{
            \makecell[{{p{\linewidth}}}]{
                \texttt{\tiny{[GM$|$GM]}}
                \texttt{game\_result = WIN} \\
            }
        }
    }
    & & \\ \\

\end{supertabular}
}

\end{document}
